\documentclass[12pt,a4paper]{article}

\usepackage[spanish]{babel}
\usepackage[utf8]{inputenc}
\usepackage[T1]{fontenc}
\usepackage{geometry}
\usepackage{amsmath}
\usepackage{booktabs}
\usepackage{hyperref}

\geometry{margin=2.5cm}

\title{Practica 3\\Algoritmos de lineas, circunferencias e intersecciones}
\author{Lovito99}
\date{27 de abril de 2026}

\begin{document}

\maketitle

\section{Objetivo}

Implementar en C++ con OpenGL moderno los algoritmos de trazado de lineas y de circunferencias indicados en la guia, y utilizarlos para:

\begin{enumerate}
	\item Dibujar una figura compuesta a partir de lineas y circunferencias.
	\item Dibujar cinco circunferencias con centros y radios aleatorios y marcar sus intersecciones.
\end{enumerate}

\section{Algoritmos implementados}

En el proyecto se implementaron los siguientes algoritmos:

\begin{itemize}
	\item Bresenham para lineas.
	\item Circunferencia por simetria de dos vias.
	\item Circunferencia por simetria de cuatro vias.
	\item Circunferencia por simetria de ocho vias.
	\item Circunferencia por punto medio.
\end{itemize}

\section{Estructura del proyecto}

\begin{center}
\begin{tabular}{ll}
\toprule
Archivo & Descripcion \\
\midrule
main.cpp & Inicializacion de OpenGL y control de escenas \\
primitivas.cpp & Algoritmos de lineas, circunferencias e intersecciones \\
escenas.cpp & Construccion de la figura y de la escena aleatoria \\
shader.vert / shader.frag & Shaders para dibujar puntos con color \\
\bottomrule
\end{tabular}
\end{center}

\section{Descripcion de la solucion}

La aplicacion genera dos escenas. La primera escena dibuja una figura compuesta mediante segmentos de recta y circunferencias construidas con distintos algoritmos. La segunda escena genera cinco circunferencias con una semilla aleatoria distinta en cada ejecucion y marca los puntos de interseccion usando cruces dibujadas con el algoritmo de lineas.

Para detectar intersecciones entre dos circunferencias se empleo la formulacion geometrica clasica a partir de la distancia entre centros. Si la distancia es compatible con una interseccion real, se calculan uno o dos puntos y despues se rasterizan marcadores sobre esos puntos.

\section{Compilacion y ejecucion}

Desde la raiz del proyecto:

\begin{verbatim}
cmake -S . -B build
cmake --build build
./build/circunferencias
\end{verbatim}

Controles del programa:

\begin{itemize}
	\item Tecla 1: muestra la figura compuesta.
	\item Tecla 2: muestra cinco circunferencias aleatorias con intersecciones.
	\item Tecla R: regenera la escena aleatoria.
\end{itemize}

\section{Conclusiones}

La practica permite comparar varios algoritmos de trazado de circunferencias basados en simetria y en decision incremental. Ademas, muestra como combinar primitivas rasterizadas para construir escenas mas complejas y como identificar geometricamente intersecciones entre circunferencias.

\end{document}
